% !TeX spellcheck = en_US
\documentclass[french]{yLectureNote}

\title{Algèbre linéaire 2}
\subtitle{Langage mathématique}
\author{Paulhenry Saux}
\date{\today}
\yLanguage{Français}

\professor{C.Dartyge}

\usepackage{graphicx}%----pour mettre des images
\usepackage[utf8]{inputenc}%---encodage
\usepackage{geometry}%---pour modifier les tailles et mettre a4paper
%\usepackage{awesomebox}%---pour les boites d'exercices, de pbq et de croquis ---d\'esactiv\'e pour les TP de PC
\usepackage{tikz}%---pour deiffner + d\'ependance de chemfig
\usepackage{tkz-tab}
\usepackage{chemfig}%---pour deiffner formules chimiques
\usepackage{chemformula}%---pour les formules chimiques en \'equation : \ch{...}
\usepackage{tabularx}%---pour dimensionner automatiquement les tableaux avec variable X
\usepackage{awesomebox}%---Pour les boites info, danger et autres
\usepackage{menukeys}%---Pour deiffner les touches de Calculatrice
\usepackage{fancyhdr}%---pour les en-t\^ete personnalis\'ees
\usepackage{blindtext}%---pour les liens
\usepackage{hyperref}%---pour les liens (\`a mettre en dernier)
\usepackage{caption}%---pour la francisation de la l\'egende table vers Tableau
\usepackage{pifont}
\usepackage{array}%---pour les tableaux
\usepackage{lipsum}
\usepackage{yFlatTable}
\usepackage{multicol}
\newcommand{\Lim}[1]{\lim\limits_{\substack{#1}}\:}
\renewcommand{\vec}{\overrightarrow}
\newcommand{\bz}{\overline{z}}
\DeclareMathOperator{\card}{card}
\renewcommand{\vec}{\overrightarrow}
\newcommand{\N}[0]{\mathbb{N}}
\newcommand{\R}[0]{\mathbb{R}}
\newcommand{\C}[0]{\mathbb{C}}
\newcommand{\dd}[0]{\mathrm{d}}
\begin{document}
\setcounter{chapter}{1}
\chapter{Applications linéaires}
\section{Application linéaire}
\begin{definition}[Application linéaire]
Soit f une application. Elle est linéaire si \(f(a+b) = f(a)+f(b)\) et \(f(\lambda a) = \lambda f (a)\), \(f(0) = 0\). Cet ensemble est noté \(\mathcal{L}\).
\end{definition}
\begin{definition}[Endomorphisme]
Un endomorphisme est une application allant dans le m\^eme espace que celui de départ.
\end{definition}
\begin{definition}[Isomorphisme]
Un endomorphisme est une application linéaire bijective de E dans E'
\end{definition}
\section{Images et noyaux}
\begin{proposition}
L'image d'un SEV par une AL est un SEV de l'espace d'arrivé
\end{proposition}
\begin{definition}[Noyau]
Noté \(\ker(f)\), c'est l'ensemble des vecteurs qui donnent le vecteur nul par l'AL.
\end{definition}
\begin{definition}[Image]
Noté \(\Im(f)\), c'est l'ensemble des vecteurs de l'espace d'arrivée que l'application peut produire.
\end{definition}
\begin{proposition}[Injectivité et noyau]
f est injective \(\iff Ker(f) = \{0\}\)
\end{proposition}
\begin{proposition}[Famille et AL]
Si f est injective et la famille E est libre, l'image de cette famille est aussi libre. Si f est surjective et la famille E est génératrice, l'image de cette famille est aussi génératrice.
\end{proposition}
\begin{definition}[Espaces isomoprphes]
2 espaces sont isomorphes \(\iff\) leur dimension est égale.
\end{definition}
\begin{theorem}[Théorème du rang]
Soit \(f:E\to E'\). ALors \[\dim(E) = rg(f)+\dim(\ker(f))\]
\end{theorem}
\begin{theorem}[Injectivité/Surjectivité en dimension finie]
Si f navigue dans 2 espaces de m\^eme dimension, les propriétés suivantes sont équivalentes : injectivité ; surjectivité ; bijectivité
\end{theorem}
\section{Matrices et applications linéaires}
\subsection{Définitions}
\begin{definition}[Matrice de f]
On appelle matrice de f dans les bases \(\{e_1,\dots,e_n\}, \{\epsilon_1,\dots, \epsilon_n\}\) \[M(f) = \begin{bmatrix}
a_{11}&a_{12}&\dots& a_{1n}\\
a_{21}&a_{22}&\dots& a_{2n}\\
\end{bmatrix}\] avec \(f(e_1) = a_{11}\epsilon_{1}+a_{21}\epsilon_{2}+\dots+a_{p1}\epsilon_p\)
\end{definition}
\subsection{Changement de base}
\begin{definition}[Matrice de passage]
Soient 2 bases de E, on peut écriree \(e_1' = p_{11}e_1+p_{21}e_{2}+\dots+p_{n1}e_n\) et la matrice de passage est  \[ \begin{bmatrix}
p_{11}&p_{12}&\dots& p_{1n}\\
p_{21}&p_{22}&\dots& p_{2n}\\
\end{bmatrix}\]
\end{definition}
\end{document}


